Foundation model-based multi-agent systems are rapidly deployed in production, yet the operational layers that ensure their reliability and efficiency remain fragmented. While orchestration frameworks and single-model inference optimization are well-surveyed, the diagnostics layer (failure localization, observability, safety) and the infrastructure layer (memory management, scheduling, resource allocation) for multi-agent workloads lack joint analysis. We present, to our knowledge, the first survey to jointly analyze both layers, covering 112 references and 45 systems across the diagnostics and infrastructure layers, together with 8 benchmarks. First, we construct a diagnostics taxonomy spanning three failure localization paradigms (statistical: FAMAS, 66\% top-1 accuracy; LLM-guided: TrajAudit, +24.7pp accuracy; RL-trained: AgenTracer, 69.10\% agent-level accuracy), three safety paradigms (topology-based, policy verification, ML-layered), observability standards, and failure benchmarks. Second, we identify three multi-agent workload characteristics (sparse activation, invocation distance, inter-agent memory sharing) that general serving systems do not exploit predictively, analyzed through application-aware systems spanning memory optimization (ScaleSim, up to 1.74$\times$ speedup), prefix-aware scheduling (TokenCake), cache lifecycle (Continuum), and power management (KAIROS). Third, we analyze the \emph{diagnostics--infrastructure tension}: diagnostics must retain runtime state while serving infrastructure must discard it, a conflict no surveyed system manages explicitly, and we derive four co-design opportunities from it. Claims used for cross-system comparison come from full-text sources; abstract-level results are marked \emph{[a]} and are never used to rank systems or compute deltas.